% ============================================================
% RESUMO — versão original de Juliane com acréscimos mínimos
% Mudanças: (1) "conduziu" -> "me levou" (regra 14); (2) frase curta
% sobre as duas linhas do capítulo 2; (3) enxertos "não só sobre, mas
% com" e auditabilidade na frase das inscrições; (4) "o conceito de
% tecnografias que proponho", direto, sem trajetória; (5) palavra-chave.
% ============================================================
\begin{center}
  \large{\textbf{Resumo}}
\end{center}
\selectlanguage{brazilian}
\begingroup
\setstretch{1.2}
\noindent O que cientistas e engenheiros fazem quando desenvolvem inteligência artificial? Movida por essa questão, segui pesquisadores ao longo de quatro anos de trabalho etnográfico no Centro de Inteligência Artificial da Universidade de São Paulo, um dos Centros de Pesquisa Aplicada da Fundação de Amparo à Pesquisa do Estado de São Paulo. O centro foi criado em 2020 mediante convênio entre a fundação, a universidade e a International Business Machines Corporation. Quando iniciei a pesquisa em 2022, a parceria parecia estável; quando a finalizei em 2025, a corporação havia encerrado o convênio. O período estendido de campo permitiu descrever o arranjo das práticas tecnocientíficas da rede sociotécnica que associou universidade pública e corporação transnacional na pesquisa de inteligência artificial no Brasil, da fundação à dissolução. Ao seguir Fábio e Cláudio, diretores do Centro, a pesquisa atravessou a história da relação entre o centro e a corporação e as práticas tecnocientíficas corporativas ao longo do tempo. Ao seguir Marcelo Finger, coordenador do projeto Sistema de Detecção Precoce de Insuficiência Respiratória baseada em Análise de Áudio, a pesquisa seguiu a cadeia de translações que converteu gravações de fala de pessoas hospitalizadas com Covid-19 em espectrogramas (representações visuais do som) processados por redes neurais para detecção de insuficiência respiratória através da análise da voz. Esse movimento etnográfico de seguir os atores por onde quer que fossem levou a tese por vários sítios, e foi esse movimento que me levou à escolha da teoria ator-rede como orientação teórico-metodológica capaz de dar conta do campo. No percurso, uma análise lexicométrica de seis obras de Bruno Latour e um mapeamento bibliométrico dos estudos sobre inteligência artificial nas ciências humanas brasileiras compõem a revisão de literatura da tese. Ao longo da pesquisa, feita não só \textit{sobre}, mas \textit{com} a inteligência artificial generativa, a produção de conhecimento etnográfico mobilizou um conjunto de práticas de inscrição tecnocientífica: diagramas interativos em Mermaid, análises em Python, modelos de linguagem (Claude, Anthropic), espectrogramas, construção de site para apresentação da defesa, com auditabilidade pública das técnicas em repositórios do GitHub. Essas grafias tecnocientíficas constituem o conceito de tecnografias que proponho, designando o modo de pesquisa que habita a tensão entre a circulabilidade técnica das inscrições e a realidade sensível dos corpos humanos. A pesquisa contribui para os estudos de ciência e tecnologia ao demonstrar que seguir múltiplos atores em um centro de pesquisa em inteligência artificial permite analisar tanto as negociações institucionais e corporativas que sustentam a pesquisa em inteligência artificial quanto as práticas tecnocientíficas situadas de desenvolvimento de sistemas, investigando como esses dois planos se articulam na produção de conhecimento em inteligência artificial e na produção de conhecimento sobre inteligência artificial pelas ciências sociais.

\noindent \textbf{Palavras-chave}: Inteligência Artificial; Tecnografias; Teoria Ator-Rede; Práticas Tecnocientíficas; Estudos de Ciência e Tecnologia.
\endgroup
\selectlanguage{brazilian}